\documentclass[11pt,a4paper]{article}
\usepackage[T1]{fontenc}
\usepackage{lmodern,amsmath,amssymb,booktabs,microtype}
\usepackage[margin=25mm]{geometry}
\usepackage[colorlinks=true,urlcolor=blue,linkcolor=blue]{hyperref}
\newcommand{\E}{\mathbb E}
\newcommand{\R}{\mathcal R}
\title{From four-leg mass filters to\\correlated time smearing in dS$_3$}
\author{Alan Day \and Chengheng\thanks{AI research assistant (OpenAI).}}
\date{Working research note --- 31 October 2026}
\begin{document}\maketitle
\begin{abstract}
The Gaussian mass regulator used in the existing static-patch trace
argument is not four independent time smearings. We decompose it exactly
into a mean-mass window and a correlated Gaussian time convolution.
For four legs the relative pair differences retain their variance,
while the pair-center variance is halved. We apply this identity to
the real-principal eikonal model, derive its small-filter correction,
and test a separately stipulated mass-averaging model.
This refines the comparison target; it neither constructs a physical
finite-$G$ state nor resolves the known trace obstruction.
\end{abstract}

\section{The source quantity and the exact filter identity}
Set $\ell=1$. CSTY (2.3) defines the mixed representation
\begin{equation}
 F(\theta_i\mid\mu)=\frac{e^{-2\pi\mu}}{Z_{\rm sphere}(\mu)}
 \int d^4m_i\,\delta\left(\mu-\frac14\sum_i m_i\right)
 e^{\sum_i(\theta_i-\theta_{i-1})m_i}\boldsymbol F(m_i).
 \label{eq:mixed}
\end{equation}
Its proposed spectral trace identification and four-leg Gaussian
regulator occur explicitly in CSTY Section 3.1 \cite{csty}.
They are prior constructions, not new states introduced here.
At equal Euclidean spacings $\pi/2$, the positive mass filter is
$s_\delta(m_i)=\exp[-\delta^2(m_i-M)^2]$.
Write $\mu=\frac14\sum_i m_i$ and $s_i=m_i-\mu$. Then exactly
\begin{equation}
 \prod_i s_\delta(m_i)
 =e^{-4\delta^2(\mu-M)^2}\,e^{-\delta^2\sum_i s_i^2}. \label{eq:split}
\end{equation}
Restricting only $\mu$ drops the second factor.

Define the cyclic difference matrix by $(Dm)_i=m_i-m_{i+1}$, with
$m_5=m_1$, and $P=I-\boldsymbol1\boldsymbol1^T/4$. If
\[
 C=\frac{\delta^2}{8}
 \begin{pmatrix}5&-1&-3&-1\\-1&5&-1&-3\\-3&-1&5&-1\\-1&-3&-1&5\end{pmatrix},
\]
then $C=2\delta^2(DD^T)^+$ and $D^TCD=2\delta^2P$.
For a centered Gaussian $d$ of covariance $C$ on $\sum d_i=0$,
\begin{equation}
 \E_d e^{i(Dm)\cdot d}=e^{-\delta^2\sum_i s_i^2}. \label{eq:char}
\end{equation}
The common-time zero mode is immaterial. This is the exact
quadratic transfer factor left as an order estimate in the source
argument. It is an application of elementary Gaussian Fourier algebra,
not a priority claim.

\newpage
\section{Correlations, normalization, and the remaining mass average}
Use the same insertion ordering as the preceding real-principal
eikonal calculation:
\[
 \theta=(-\pi/2+i(t+d_1),\,id_2,\,\pi/2+i(t+d_3),\,\pi+id_4).
\]
The combinations
\[
 X=(d_1-d_3)/2,\qquad Y=(d_2-d_4)/2,\qquad
 Z=(d_1+d_3-d_2-d_4)/2
\]
are independent with
\begin{equation}
 \operatorname{Var}X=\operatorname{Var}Y=\operatorname{Var}Z
 =\delta^2/2. \label{eq:variance}
\end{equation}
For four independent insertion shifts of variance $\delta^2$,
the first two variances agree but $\operatorname{Var}Z=\delta^2$.
Thus the two protocols cannot be interchanged.

Denote this correlated convolution by $\mathcal C_\delta$.
Inserting \eqref{eq:char} into the spectral integral gives the
unnormalized filtered quantity, under the same spectral trace
hypothesis as the source:
\begin{equation}
 {\cal T}_{M,\delta}=
 \int d\mu\, e^{2\pi\mu} Z_{\rm sphere}(\mu)
 e^{-4\delta^2(\mu-M)^2}
 \mathcal C_\delta F(\theta_i\mid\mu). \label{eq:average}
\end{equation}
All spectral support restrictions must be retained in $F$.
Equation \eqref{eq:average} is also a formal transform identity if
that hypothesis is not assumed. In that case it is not established
as a positive trace or a state preparation.

For the specified eikonal model, take real principal labels
$\nu,\nu'>0$, probe masses squared $1+\nu^2,1+\nu'^2$,
$\gamma(\mu)=4G\mu$, and $\alpha(\mu)=2\pi(1-\gamma(\mu))$.
Let
\[
 I_{n,\nu}(X)=\int_0^\infty x^{n+1}
 K_{i\nu}(xe^X)K_{i\nu}(xe^{-X})\,dx,\qquad
 B_{n,\nu}(\delta)=\E_X I_{n,\nu}(X).
\]
The disconnected term is still
$D_\delta(\mu)=B_{0,\nu}B_{0,\nu'}/[\pi\alpha(\mu)]^2$.
If the numerator in \eqref{eq:average} is divided by the corresponding
disconnected integral, its normalized version is
\begin{equation}
 \overline\R=
 \frac{\int d\mu\,a(\mu)D_\delta(\mu)\R^{\rm corr}_\delta(t\mid\mu)}
 {\int d\mu\,a(\mu)D_\delta(\mu)},\qquad
 a(\mu)=e^{2\pi\mu}Z_{\rm sphere}(\mu)e^{-4\delta^2(\mu-M)^2}.
 \label{eq:ratio}
\end{equation}
It is a positive average only if $aD_\delta$ is positive and integrable.
The source normalization constant and actual time-ordered
factorization remain separate questions. In particular, \eqref{eq:ratio}
does not identify an arbitrary $G=0$ observer wavepacket with the
mixed representation. At fixed theory $G$ stays fixed during the
$\mu$ integral; holding $\gamma$ fixed would average different theories.

\newpage
\section{A calculable change in the real-principal coefficients}
The previous lateral kernel representation remains valid with
\eqref{eq:variance}. In detail, put
\[
 A=e^Y\cosh r_1+e^{-Y}\cosh r_2,\quad
 B=e^X\cosh r_3+e^{-X}\cosh r_4,\quad T=t+Z,
\]
and $W=\cos\nu r_1\cos\nu r_2\cos\nu'r_3\cos\nu'r_4$.
For $\eta=+1$ the continued kernel is
\[
 {\cal K}_\eta(c)=-J_\eta'(c),\quad
 J_\eta(-1/a)=ae^{-a}\operatorname{Ei}(a)-\eta i\pi ae^{-a}.
\]
The conditional normalized model is
\[
 \R^{\rm corr}_\eta=
 \frac{\E\left\langle\int_{\mathbb R_+^4}
 \frac{W}{A^2B^2}
 \left[{\cal K}_\eta(he^T/(AB))+{\cal K}_{-\eta}(he^{-T}/(AB))-1\right]
 d^4r\right\rangle_\phi}{B_{0,\nu}B_{0,\nu'}}.
\]
Here $h=-\cos[(1-\gamma)(\pi-\phi)]/[2\mu\,\operatorname{sinc}(\pi\gamma)]$.
It retains both phases. At each fixed perturbative order the factor is
\begin{equation}
 \frac{(-1)^n\langle h^n\rangle}{n!}
 e^{n^2\delta^2/4}\frac{B_{n,\nu}B_{n,\nu'}}{B_{0,\nu}B_{0,\nu'}}
 (e^{nt}+e^{-nt}). \label{eq:coeff}
\end{equation}
The independent-time prescription instead has $e^{n^2\delta^2/2}$.
Moment positivity still follows from its established Mellin spectrum;
it does not prove positivity of the full gravitational theory.

The small-width change is explicit. With $a=(n+2)/2$ and
$f_n(r)=e^{ar}K_{i\nu}(e^r)$, integration by parts in the Bessel
equation gives
\[
 \frac{\|f_n'\|_2^2}{\|f_n\|_2^2}
 =a^2+\nu^2-\frac{I_{n+2,\nu}(0)}{I_{n,\nu}(0)}
 =q_{n,\nu}:=\frac{(n+2)^2/4+\nu^2}{n+3}.
\]
Boundary terms vanish for $n\geq0$ and real $\nu$.
The Mellin transform and Gaussian multiplier then imply
\[
 B_{n,\nu}(\delta)=I_{n,\nu}(0)
 [1-q_{n,\nu}\delta^2+O(\delta^4)].
\]
The coefficient in \eqref{eq:coeff}, relative to zero width, changes by
$1+c_n\delta^2+O(\delta^4)$, where
\[
 c_n=\frac{n^2}{4}-q_{n,\nu}+q_{0,\nu}
                  -q_{n,\nu'}+q_{0,\nu'},\qquad
 c_1=\frac{2(\nu^2+\nu'^2)-5}{24}.
\]
For $(\nu,\nu')=(.6,.9)$, $c_1=-.110833\ldots$; independent
time smearing gives $+.139167\ldots$. The sign of this width
correction differs, while the first gravitational growing term itself
still has the positive sign of the existing eikonal model.

\newpage
\section{Endpoint control and an explicitly stipulated comparison}
The mass integral is not harmless merely because the Gaussian is
small far from $M$. Combining the large-mass approximation
$Z_{\rm sphere}(\mu)\propto\mu^2e^{-2\pi\mu}$ with
$D_\delta\propto\alpha^{-2}$ over $0<\mu<\kappa=1/(4G)$
would give a weight proportional to
\[
 \frac{\mu^2e^{-4\delta^2(\mu-M)^2}}{(\kappa-\mu)^2},
\]
whose integral diverges at $\mu=\kappa$ for every finite $\delta$.
This is failure of an extrapolated approximation, not a physical
prediction of divergent trace normalization.

For a concrete finite comparison, we separately stipulate the
semiclassical relation $\alpha=2\pi(1-\mu/\kappa)$ at finite $\kappa$
and use the CLP proposal
$Z_{\rm obs}(\alpha)=32\sin^2(\alpha/2)\sinh^2(\kappa\alpha/2)$
\cite{clp}. This finite-$\kappa$ map and the insertion completion
are not established. After removing a common positive constant,
the weight in \eqref{eq:ratio} is exactly
\begin{equation}
 A_\kappa(\mu)e^{-4\delta^2(\mu-M)^2},\qquad
 A_\kappa(\mu)=
 \left[\frac{\sin(\alpha/2)}{\alpha}
 (1-e^{-\kappa\alpha})\right]^2,\quad0<\mu<\kappa.
 \label{eq:A}
\end{equation}
It satisfies $0\leq A_\kappa\leq1/4$ and tends to
$\pi^2(\kappa-\mu)^2$ at the upper endpoint.
The mass average of the bounded conditional eikonal ratio therefore
exists. This restores an integrable comparison weight, not an
established physical finite-$G$ density matrix.

On any fixed interior relative mass window, the preceding
rapidity representation has bounded second derivatives in
$r=\mu/M$ and in the three displacement variables, for fixed
principal labels, $\gamma(M)$ away from one and $t=\log M+u$
with $u$ in a compact set. The bound follows by differentiating
the smooth real lateral kernel; its derivatives are bounded,
Gaussian exponential moments are finite, and the rapidity
majorant is integrable. Taylor expansion with centered Gaussians
then bounds the conditional smearing error by $O(\delta^2)$
and the normalized mean-mass averaging error by
$O((\delta^2 M^2)^{-1})$, provided the positive density has bounded
relative $C^2$ derivatives there. For \eqref{eq:A}, its endpoint
bound supplies an exponentially small exterior probability when
$\delta^2\asymp M^{-1}$ and $\gamma(M)$ is fixed in $(0,1)$.
Thus the source's $O(M^{-1})$ balancing estimate holds for this
stipulated comparison model. The constants are not bounds on
omitted gravitational terms.

\paragraph{What is still missing.}
CLP (3.35)--(3.37) only calibrates the physical insertion in a
leading heavy-heavy-light limit. Its positive two-point expression
does not fix the finite-parameter four-point Gram inverse or
Lorentzian state. Our filter identity neither supplies these data
nor invalidates the source's positivity argument, which already
uses the four separate mass filters.

\newpage
\section{Numerical checks and reproducibility}
Exact rational algebra checks $D^TCD=2\delta^2P$ and the reduced
covariance. A three-level Hermitian matrix example compares the
actual filtered trace with independent Gaussian quadrature of the
loop transform; both agree to the saved tolerance. The mean-only
window is also retained and gives a different result. Mellin
second moments independently test $q_{n,\nu}$ at 50 decimal digits.
All 23 diagnostics pass; this finite matrix check is not a test of
quantum-gravity dynamics.

For the conditional eikonal integrals we use $\nu=.6,\nu'=.9$,
$\mu=1000$, $\gamma=.25$, $t=\log1000+u$. Eight scrambled Sobol
seeds at $2^{19}$ and $2^{21}$ points retain both protocols with
common samples. The following values are correlated minus
independent, with empirical component standard errors across
scrambles:
\begin{center}\begin{tabular}{cc rr rr}
\toprule
$\delta$&$u$&$\Re\Delta\R$&$\Im\Delta\R$&SE$_{\Re}$&SE$_{\Im}$\\
\midrule
$1/\sqrt{1000}$&0&$-.00015810$&$.00022730$&$5.14\,10^{-8}$&$4.86\,10^{-8}$\\
.4&0&$-.007832$&$.033554$&$.0000071$&$.0000067$\\
.8&0&$.032164$&$.079138$&$.0000189$&$.0000315$\\
.8&2&$-.079416$&$-.147868$&$.0000434$&$.0000599$\\
\bottomrule
\end{tabular}\end{center}
These are full conditional kernel integrals, not just perturbative
coefficients. They are not mass-integrated traces.

A separate nine-dimensional run implements \eqref{eq:A} at
$M=1000$, $\kappa=4000$, $G=1/16000$,
$\delta=1/\sqrt{1000}$ and $u=0$, with antithetic mass offsets.
Its complete raw estimates, mass normalizer and paired differences
are retained in \path{results/stage4/mass-filter/packet.json}.
The comparison is explicitly labeled as the stipulated CLP-weight
model; a change of mass never changes $G$.
Empirical scramble errors are not rigorous total error certificates.

All drivers refuse to overwrite raw evidence. The stage records
dependency versions, code/input hashes, finite-matrix controls,
and both sampling levels. The next task is to determine a physical
finite-parameter insertion and gluing prescription, or to establish
a concrete obstruction to one, retaining the exact filter protocol
derived here.

\begin{thebibliography}{9}
\bibitem{csty} Y. Chen, D. Stanford, H. Tang and Z. Yang,
\emph{Negative shocks versus static patch holography},
\href{https://arxiv.org/abs/2607.14042v1}{arXiv:2607.14042v1},
(2.3), Section 3.1 and Appendix F.
\bibitem{clp} J. Chandra, D. Liska and B. Post,
\emph{3D de Sitter from complex Virasoro TQFT: Observers and knots},
\href{https://arxiv.org/abs/2610.30247v1}{arXiv:2610.30247v1},
(3.5), (3.10), (3.35)--(3.40).
\end{thebibliography}
\end{document}
